\section{Conclusion}\label{section::conclusion}

This thesis investigated the impact of continual multilingual pre-training on the cross-lingual transferability and task-specific performance of Large Language Models (LLMs) for text-to-SPARQL query generation, with a specific focus on the source languages English and German. 
By comparing the Occiglot-7b-eu5 model, a result of continual pre-training of Mistral 7B v0.1 on five European languages \cite{avramidis-etal-2024-occiglot}, against the original Mistral model, this work shows the nuanced effects of such pre-training strategies.

A primary contribution of this research was the development and public release of a substantial, reproducible multilingual dataset for text-to-SPARQL generation, comprising 895,954 examples across 12 languages. 
This dataset, derived from established benchmarks like QALD \cite{usbeck2024qald} and other knowledge extraction challenges \cite{jiang2022knowledge}, and augmented with German translations and Wikidata entity/relationship mappings, facilitated the fine-tuning and evaluation of eight distinct model variants. 
The study also contributed these fine-tuned models and an custom evaluation framework to the research community.

The experimental findings reveal a complex interplay between multilingual pre-training, task characteristics, and the availability of contextual information. 

Key conclusions drawn are:

\begin{itemize}
    \item Continual multilingual pre-training, as embodied by Occiglot, demonstrably enhances performance in newly introduced languages (German) when explicit entity/relationship mappings are absent (v1 models). 
    This underscores the value of targeted pre-training for improving lexical and semantic understanding in specific languages.
    \item The provision of explicit entity/relationship ID mappings (v2 models) substantially mitigates the performance differences attributable to language-specific pre-training. 
    This suggests that contextual grounding can compensate for a lack of extensive target-language pre-training.
    The entity/relationship linking was identified as the primary bottleneck for v1 models, with v2 models showing approximately fourfold performance increases when mappings were supplied.
    \item The impact of Occiglot's multilingual pre-training on its English performance is nuanced. 
    While Occiglot maintained comparable or sometimes superior executable query generation rates in English, Mistral significantly outperformed Occiglot in F1-score on the complex QALD-10 English benchmark under both fine-tuning conditions (v1 and v2).
    This points to a potential trade-off where multilingual adaptation might affect nuanced handling of the original language, an observation that finds resonance with reports from the Occiglot research team \cite{avramidis-etal-2024-occiglot, OcciglotResearchCollective2024Occiglot7BEU5Instruct}.
    \item The nature of the test dataset influences observed performance. 
    The QALD-10 benchmark, with its inherent query complexity \cite{usbeck2024qald}, posed challenges distinct from the other test sets. 
    Nevertheless, competitive F1-scores (around 0.25) were achieved on QALD-10 with v2 models, comparable to other systems in the QALD-10 benchmark \cite{usbeck2024qald}.
\end{itemize}

This research frames text-to-SPARQL generation as a structured output task \cite{goodfellow2016deep}, and the findings contribute to understanding the conditions under which multilingual pre-training is most beneficial. 
While this study faced limitations, including its focus on high-resource languages and the reliance on LLMs for certain data processing steps (as detailed in Section 8.2), it extends the work of Rangel et al. \cite{rangel2024sparqlgenerationanalysisfinetuning} by exploring the dynamics of the task in a multilingual setting.

Ultimately, this work underscores that enhancing LLM performance across diverse languages involves more than just increasing multilingual data; it requires careful consideration of task complexity, linguistic challenges, and potential capability trade-offs. 
The datasets, models, and insights presented herein aim to provide a foundation for future research directed at creating more equitable and effective LLMs, thereby helping to bridge the linguistic divide in access to information and AI-driven efficiencies.